\paragraph{Constructing both filtrations from the biextension.}
The two filtrations can be built on the original family itself.
The construction uses first infinitesimal neighbourhoods and
multiplication by $p$.

\begin{theorem}\label{geo-biextension-conjugate-filtration}
Let $p$ be a prime, let $S$ be a scheme over $\mathbf F_p$,
and let $G/S$ be a
smooth separated commutative group scheme of relative dimension $d$.
Let $\mathcal P$ be a symmetric biextension of $G\times_S G$ by
$\mathbf G_m$, with its rigidifications on the axes and its symmetry
equal to the identity on the diagonal.
Suppose that, at every geometric point of $S$, the group is
semiabelian and its biextension is the pullback from the Poincar\'e
biextension of a line bundle inducing an isomorphism between its
abelian quotient and the dual quotient.

Put $\mathcal V=\operatorname{Lie}(G/S)$ and
$\mathcal U=\mathcal V^*$.
There is a vector bundle $\mathcal H$ on this same base with exact
sequences
\begin{equation}\label{geo-actual-two-filtrations}
 0\longrightarrow\mathcal U\longrightarrow\mathcal H
       \longrightarrow\mathcal V\longrightarrow0,\qquad
 0\longrightarrow F_S^*\mathcal V\longrightarrow\mathcal H
       \longrightarrow F_S^*\mathcal U\longrightarrow0.
\end{equation}
It has the formal-coordinate transition functions and integral
lattices defined by the first derivatives and mixed bilinear
coefficients of $\mathcal P$.
An action of a commutative coefficient ring with involution
$a\mapsto\bar a$ on $G$, with compatible adjunction on $\mathcal P$,
makes both sequences equivariant; the action on $\mathcal U$ is
the conjugate dual action.
\end{theorem}
\begin{proof}
Let $G_1$ be the first infinitesimal neighbourhood of the zero
section. It is finite locally free over $S$ of rank $d+1$, with
square-zero ideal $\mathcal U$.
Over a point $x$ of $G$, consider pointed sections of
$\mathcal P_x|_{G_1}$. They form a torsor under the additive vector
group $\mathcal U$.
These torsors are representable: push the line bundle on
$G\times_S G_1$ forward to $G$, and take the inverse image of the
unit section under restriction to $G\times_S S$.
The first biextension law gives this torsor a commutative group
law. Denote the resulting vector extension by $\mathcal E$:
\[
        0\longrightarrow\mathcal U_{\mathrm{add}}
          \longrightarrow\mathcal E\xrightarrow{\pi}G
          \longrightarrow0.
\]
Taking Lie algebras gives the first sequence, with
$\mathcal H=\operatorname{Lie}(\mathcal E/S)$.
All these constructions commute with base change.

Multiplication by $p$ on $\mathcal E$ kills its additive kernel.
It therefore descends to a group homomorphism $j:G\to\mathcal E$
with $j\pi=[p]_{\mathcal E}$.
Differentiation at the identity gives $dj=0$, since $d\pi$ is
surjective and $d[p]_{\mathcal E}=0$. Translation gives zero
relative differential everywhere.
For a smooth scheme in characteristic $p$, functions with zero
relative differential are precisely the image under relative
Frobenius; this follows in polynomial coordinates and passes to
\'etale coordinates. Hence there is a unique homomorphism
\[
       \Phi:G^{(p)}\longrightarrow\mathcal E,\qquad
       \Phi F_{G/S}=j.
\]
Its differential is a map
$\alpha:F_S^*\mathcal V\to\mathcal H$.

We first verify its exact coordinate formula and isotropy.
Choose coordinates $x$ on the formal group and a frame $e(x,y)$
of $\mathcal P$, rigidified on both axes. A pointed section on
$G_1$ is written $e(x,y)(1+\lambda(y))$.
Write $C(x,y)$ and $D(x,y)$ for the scalars of $p$-fold
biextension multiplication in the first and second variable,
respectively, and write the symmetry as
\[
               e(x,y)\longmapsto s(x,y)^{-1}e(y,x).
\]
Put
\[
 h_{ij}=[x_j^p]\,[p]_G(x)_i,\qquad
 A_{ij}=[x_i^py_j]\,C(x,y),\qquad
 \Sigma_{ij}=[x_i y_j]\,s(x,y).
\]
Multiplication in $\mathcal E$ adds the derivative in $y$ of the
first biextension cocycle to the vector coordinate. Thus
\[
 j(x)=\bigl([p]_G(x),d_yC(x,0)\bigr),\qquad
                         \alpha(v)=(hv,A^tv).
\]
In particular $dj=0$ gives $[x_i y_j]C=0$.
The biextension symmetry and interchange laws give
\begin{align*}
 D(x,y)&=C(y,x)s(x,[p]y)s(x,y)^{-p},\\
 C(x,[p]y)D(x,y)^p&=D([p]x,y)C(x,y)^p.
\end{align*}
The first identity also gives $[x_i y_j]D=0$.
Since $[p]_G$ factors through Frobenius, its terms have degree at
least $p$. Taking the coefficient of $x_i^py_j^p$ in the second
identity, and then using the first, yields
\begin{equation}\label{geo-biextension-isotropy-identity}
                  Ah=h^tA^t+h^t\Sigma h.
\end{equation}
Indeed the $p$-th powers of $C,D$ contribute zero to that
coefficient, while $s([p]x,y)^{-p}$ has first-variable degree at
least $p^2$. The three remaining contributions are exactly those
displayed in \eqref{geo-biextension-isotropy-identity}.

The symmetry of $\mathcal P$ defines a perfect alternating
pairing on $\mathcal H$. Intrinsically, compare the two pointed
sections at a pair of first-order tangent points after swapping
the two factors. In the chosen frame its value is
\begin{equation}\label{geo-biextension-pairing}
 \langle(v,\lambda),(w,\mu)\rangle
              =\mu(v)-\lambda(w)+v^t\Sigma w.
\end{equation}
The diagonal condition gives $s(x,x)=1$ and makes $\Sigma$
alternating, also when $p=2$. This description proves that the pairing is independent
of the frame and is perfect. Equation
\eqref{geo-biextension-isotropy-identity} says precisely that
$\operatorname{im}\alpha$ is isotropic.

It remains to prove injectivity of $\alpha$ on geometric fibres.
For a proper abelian fibre $B$, the extension just constructed is
its universal vector extension, after the prescribed self-duality.
The universal extension of $V_B:B^{(p)}\to B$ kills its additive
kernel: on that kernel the map is pullback by the dual
homomorphism $F_{B^\vee}$, whose differential is zero.
It consequently factors through a map $B^{(p)}\to\mathcal E$.
Functoriality and $V_BF_B=[p]_B$ identify this map with $\Phi$.

To fix the Frobenius convention, take the contravariant
Dieudonn\'e module of $B^\vee[p^\infty]$.
Its reduction is the Lie algebra of the universal extension of
$B$, and the reduction of crystalline Frobenius is the differential
of the universal extension of
$(F_{B^\vee})^\vee=V_B$.
Its kernel is the Frobenius pullback of the additive Hodge
subspace. Therefore its induced map on $F_S^*\operatorname{Lie}B$
is injective, as required
\cite[II, \S\S9 and 13]{mazur-messing-universal}
\cite[\S1.1.4, equation (1.1.4.1)]{madapusi-pera-realizations}.

For a semiabelian geometric fibre let
$0\to T\to G\xrightarrow{q}B\to0$ be its torus quotient.
The pullback biextension is canonically trivial on $G\times T$.
Restricting pointed sections to $T_1$ gives a homomorphism
$\mathcal E\to\omega_T$. Its kernel is
$\mathcal E'=G\times_B E_B$, where $E_B$ is the extension for $B$.
Since $p$ kills $\omega_T$, the map $\Phi$ lands in $\mathcal E'$.
The composite to $E_B$ is $\Phi_Bq^{(p)}$; this follows by
composing with relative Frobenius and using multiplication by $p$.
On $T^{(p)}$, the map is the toric Verschiebung into
$T\subset\mathcal E'$, which is an isomorphism over the
algebraically closed field.
If $\alpha(v)=0$, its image in the proper quotient gives
$dq^{(p)}(v)=0$. Thus $v$ lies in $\operatorname{Lie}T^{(p)}$,
where toric Verschiebung is injective. Hence $v=0$.

The fibrewise rank criterion now makes $\alpha$ a locally split
injection of vector bundles over $S$.
Define $\beta:\mathcal H\to(F_S^*\mathcal V)^*$ by
$\beta(z)(v)=\langle\alpha(v),z\rangle$.
Perfection of the pairing makes $\beta$ surjective; isotropy and
the ranks give $\ker\beta=\operatorname{im}\alpha$.
This is the second sequence in
\eqref{geo-actual-two-filtrations}.

Finally, an adjunction of the coefficient action induces
endomorphisms of $\mathcal E$. They commute with multiplication
by $p$, so the uniqueness of $\Phi$ proves equivariance of
$\alpha$ with the pulled-back action. In coordinates, compatibility
of adjunction with symmetry says
\[
 u_a(x,y)u_{\bar a}(y,x)
                    =s(x,\bar a y)/s(ax,y).
\]
Its mixed coefficient is
$b_a+b_{\bar a}^t=\Sigma\rho(\bar a)-\rho(a)^t\Sigma$.
Together with \eqref{geo-formal-hodge-action}, this proves
$\langle az,z'\rangle=\langle z,\bar a z'\rangle$.
It gives the conjugate dual action on the quotient.
The construction uses the original infinitesimal
neighbourhood and line bundle at every stalk, so its lattices are
the original integral lattices.
\end{proof}

For clarity, the coordinate formulas also specify the comparison
with the initial formal construction. Write the two frames as
$e_2(x,y)=\ell(x,y)e_1(\phi x,\phi y)$, with
$L=d\phi$ and $Z$ the mixed matrix of $\ell$. Then the transition
$T:\mathcal H_2\to\mathcal H_1$ satisfies
\begin{align*}
        T(v,\lambda)&=\bigl(Lv,L^{-t}(\lambda+Z^tv)\bigr),\\
        \Sigma_2&=L^t\Sigma_1L+Z^t-Z.
\end{align*}
The coefficient action is \eqref{geo-formal-hodge-action},
$\alpha=(h,A^t)$, and
\[
                  \beta(v,\lambda)
                        =h^t\lambda-(A-h^t\Sigma)v.
\]
The matrix $\Sigma$ must be retained, including when $Z$ is
nonsymmetric.
These formulas preserve the projection to the actual Lie lattice
and the actual dual sublattice.

\begin{corollary}[Compatible formal charts on a curve]
\label{geo-biextension-formal-curve}
The same conclusion holds on a smooth curve in characteristic $p$
when one is given
compatible formal biextension charts at its generic point and
over its local rings, with the displayed integral transition
matrices, and each geometric residue chart comes from a
semiabelian variety satisfying the theorem's fibre hypothesis.
Here no semiabelian scheme over the whole curve is required.
\end{corollary}
\begin{proof}
The first-neighbourhood construction and its multiplication law
make sense formally. Over a ring of characteristic $p$, a formal
power series with zero partial derivatives involves only powers
divisible by $p$, so the same relative Frobenius factorization
defines $\alpha$. Its formula, the pairing, and $\beta$ commute
with all the displayed transitions. They therefore give maps of
the prescribed integral vector bundles. The fibrewise argument
proves injectivity of $\alpha$ at each geometric residue point.
The locally split rank criterion, followed by isotropy, gives both
exact sequences. This also proves their coefficient equivariance.
\end{proof}

In applying this result, the line bundle itself must be tracked.
If the original line $L$ on a geometric fibre is isomorphic to
$q^*M$ and the translation homomorphism
$\lambda_M:B\to B^\vee$ has trivial scheme-theoretic kernel, then
$\lambda_M$ is an isomorphism: its kernel is finite, so its image
has the full dimension of $B^\vee$, and its degree is one.
The associated biextension is the pullback of that of $M$.
Indeed, after normalizing the line isomorphism at zero, any
discrepancy in the multiplication laws is a unit on a product
of connected reduced varieties, trivial on the coordinate axes,
and is therefore one by \cite[Theorem~1.1]{conrad-units-product}.
The canonical exchange of the two factors is identity on the
diagonal. This supplies the required symmetry also in
characteristic two. Tensoring $L$ by an untracked positive power
would not justify this argument.

Together with coefficient freeness, the result supplies the two
filtered bundles from the prescribed compatible charts and
geometric fibres. Constructing those data on the chosen integral
model, and evaluating its integral complex multiplication
determinant, remain separate inputs.
